% =============================================================================
% Project : THE_PHILOSOPHY_OF_VIABILITY
% Book    : Book I — Ontology of Viability
% Book ID : VS.B01
%
% File    : 4_ASSETS/diagrams/
%           VS.B01.P01.C01.S06.SS13.P001.DG01.identity_individuation_conceptual_relations.tex
%
% Asset:
%   Conceptual Relations of Identity, Unity, and Individuation
%
% Purpose:
%   Provides a schematic representation of the analytical questions
%   concerning an existent, distinguishing identity, unity,
%   individuation, distinction, and numerical identity without
%   implying causal, temporal, or hierarchical priority.
%
% Requires (preamble):
%   \usepackage{tikz}
%   \usetikzlibrary{positioning, fit, backgrounds}
% =============================================================================

\begin{figure}[htbp]
    \centering
    \hyphenpenalty=10000 \exhyphenpenalty=10000\relax  % no hyphenation inside diagram labels

    \begin{tikzpicture}[
        font=\small,
        node distance=0.45cm,
        % --- base box ---
        box/.style={
            draw, rounded corners=3pt, align=center,
            minimum height=1.35cm, inner sep=5pt, line width=0.6pt
        },
        % --- primary questions (three stacked, equal width) ---
        primary/.style={box, text width=9.2cm, fill=blue!8, draw=blue!55!black},
        % --- related questions (2 x 2 grid, equal width) ---
        related/.style={box, text width=4.4cm, fill=gray!10, draw=black!60},
        % --- group frames ---
        group/.style={
            draw=black!45, dashed, rounded corners=5pt,
            inner sep=9pt, line width=0.5pt
        },
        % --- captions and title (names avoid clashes with built-in TikZ keys) ---
        grouplabel/.style={font=\footnotesize\itshape, text=black!75, align=center},
        figtitle/.style={font=\small\bfseries, align=center}
    ]

    % ---------------------------------------------------------------
    % Group A: primary analytical questions
    % ---------------------------------------------------------------
    \node[primary] (identity) {
        \textbf{Identity}\\[2pt]
        The identity of the existent
    };
    \node[primary, below=of identity] (unity) {
        \textbf{Unity}\\[2pt]
        The constitution of a unity
    };
    \node[primary, below=of unity] (individuation) {
        \textbf{Individuation}\\[2pt]
        The constitution of the existent as an individual
    };

    \node[group, fit=(identity)(unity)(individuation)] (groupA) {};
    \node[grouplabel, below=0.15cm of groupA] (capA) {
        Analytically distinguishable questions
    };

    % ---------------------------------------------------------------
    % Group B: related questions (2 x 2)
    % ---------------------------------------------------------------
    \node[related, below=1.1cm of capA, xshift=-2.4cm] (distinction) {
        \textbf{Distinction}\\[2pt]
        Difference and non-identity
    };
    \node[related, right=0.4cm of distinction] (numerical) {
        \textbf{Numerical Identity}\\[2pt]
        Sameness across comparison
    };
    \node[related, below=0.4cm of distinction] (singularity) {
        \textbf{Singularity}\\[2pt]
        Singular status within a respect
    };
    \node[related, below=0.4cm of numerical] (uniqueness) {
        \textbf{Uniqueness}\\[2pt]
        Non-repetition within a comparison
    };

    \node[group, fit=(distinction)(numerical)(singularity)(uniqueness)] (groupB) {};
    \node[grouplabel, below=0.15cm of groupB] (capB) {
        Related questions; no automatic derivation or reduction
    };

    % ---------------------------------------------------------------
    % Outer frame and title
    % ---------------------------------------------------------------
    \begin{scope}[on background layer]
        \node[draw=black, rounded corners=6pt, line width=0.8pt,
              inner sep=12pt, fit=(groupA)(capA)(groupB)(capB)] (outer) {};
    \end{scope}
    \node[figtitle, above=0.25cm of outer] {
        ONE EXISTENT UNDER ONTOLOGICAL EXAMINATION
    };

    \end{tikzpicture}

    \caption{Analytical Relations Among Identity, Unity, and Individuation}
    \label{fig:identity-individuation-conceptual-relations}
\end{figure}

% =============================================================================
% END OF DIAGRAM DG01
% =============================================================================
